
\documentclass[11pt]{article}
\usepackage[margin=1in]{geometry}
\usepackage{amsmath,amssymb}
\usepackage{enumitem}
\usepackage[T1]{fontenc}
\usepackage{lmodern}
\setlength{\parskip}{0.5em}
\setlength{\parindent}{0pt}

\begin{document}

\begin{center}
{\LARGE \textbf{IB Physics 2016--2020}}\\
{\large \textbf{Mapped Collection for E1}}\\
\vspace{0.3em}
\textit{Near-literal paraphrases matched to atomic models, spectra, photoelectric effect, and matter waves.}
\end{center}

\textbf{Use.} These are intentionally close in logic and context to the original questions so that working through them feels much like doing the paper itself.

\section*{Problems}

\begin{enumerate}[leftmargin=*, itemsep=1.1em]

\item \textbf{M17 TZ2 HL Paper 2, Q5(d) --- helium identified from emission lines}\\
\textit{Context.} Rutherford and Royds collected the gas formed when alpha particles were stopped in matter and identified that gas by looking at its line spectrum.
\begin{enumerate}[label=(\alph*)]
\item Explain how an emission spectrum is produced in terms of atomic energy levels and photon emission.
\item State why the observed spectrum can be used to identify helium.
\end{enumerate}

\item \textbf{M18 TZ2 HL Paper 2, Q9 --- from Rutherford to Bohr, then to hydrogen orbits}\\
\textit{Context.} A question traces the development of the atomic model from Rutherford's scattering results to Bohr's quantized orbits and then applies the model to hydrogen.
\begin{enumerate}[label=(\alph*)]
\item Describe Rutherford's model of the atom.
\item State what features of alpha-scattering observations support that model.
\item Bohr proposed the condition
\[
mvr=\frac{nh}{2\pi}.
\]
Explain why this modification was introduced.
\item For a hydrogen atom, show that the electron speed $v$ in a circular orbit of radius $r$ satisfies the required relation by equating electrostatic attraction to centripetal force.
\item For the state with quantum number $n=6$, calculate the electron speed.
\item Hence determine the corresponding orbital radius.
\end{enumerate}

\item \textbf{M19 TZ1 HL Paper 2, Q2 --- electron scattering and nuclear size}\\
\textit{Diagram.} A beam of electrons, each with de Broglie wavelength $2.4\times10^{-15}\,\text{m}$, is directed at a thin film of silicon-30 nuclei. A graph is given of scattered-electron intensity against scattering angle. The graph indicates a first minimum from which the nuclear diameter can be estimated.
\begin{enumerate}[label=(\alph*)]
\item Use the graph to show that the radius of a silicon-30 nucleus is of the order of $4\,\text{fm}$.
\item Use your result and an appropriate nuclear-radius scaling law to estimate the radius of a thorium-232 nucleus.
\item Suggest one reason why electrons of a given energy are more suitable than alpha particles of the same energy for investigating nuclear size.
\end{enumerate}

\item \textbf{N19 HL Paper 2, Q8 --- singly ionized helium: classical model, Bohr model, and Schr\"odinger model}\\
\textit{Diagram.} A singly ionized helium atom is modelled first as a nucleus with charge $+2e$ and one electron moving in a circular orbit of radius $r$. A later part states that, in the Bohr model for this ion, the orbital radius is
\[
r = 2.7\times10^{-11}n^{2}\,\text{m}.
\]
\begin{enumerate}[label=(\alph*)]
\item In the classical circular-orbit model, show that the electron speed is
\[
v=\sqrt{\frac{2ke^{2}}{mr}}.
\]
\item Hence deduce that the total energy of the electron is
\[
E_{\text{TOT}}=-\frac{ke^{2}}{r}.
\]
\item If the orbiting electron continually loses energy by radiation, describe what happens to the orbital radius according to the classical model.
\item For the Bohr model with $n=3$, show that the electron de Broglie wavelength is about $5.1\times10^{-10}\,\text{m}$.
\item Estimate, to one significant figure, the ratio \emph{orbit circumference / electron wavelength}.
\item Compare the Bohr model with the Schr\"odinger model of the atom.
\end{enumerate}

\item \textbf{N19 HL Paper 2, Q11 --- photoelectric effect}\\
\textit{Diagrams.} One part refers to monochromatic light of very low intensity incident on a metal, with electrons emitted almost immediately. Another part shows a photoelectric-effect apparatus and a graph of photocurrent $I$ against cathode potential $V$ for light of wavelength $480\,\text{nm}$.
\begin{enumerate}[label=(\alph*)]
\item Explain why the immediate emission of electrons does \emph{not} support a classical wave model of light.
\item Explain why the same observation \emph{does} support the photon model.
\item Use the graph and wavelength data to determine the work function of the metal, in eV.
\item Sketch how the $I$-$V$ graph changes if the light intensity is halved while the wavelength is unchanged.
\end{enumerate}

\item \textbf{N20 HL Paper 2, Q10 --- electron diffraction from nuclei}\\
\textit{Diagram.} A beam of electrons is used to probe carbon nuclei. Data are provided for the electron momentum or energy, and a graph of scattered intensity against angle indicates a diffraction minimum.
\begin{enumerate}[label=(\alph*)]
\item Determine the electron de Broglie wavelength.
\item Explain why the scattering pattern implies diffraction and therefore wave behaviour.
\item Use the diffraction geometry to calculate the scattering angle associated with the minimum.
\item Explain why electrons with a much longer de Broglie wavelength are unsuitable for resolving nuclear dimensions.
\item Show that nuclear density is approximately independent of nucleon number $A$.
\end{enumerate}

\end{enumerate}

\end{document}
